\documentclass{beamer}

\usepackage{amssymb}
\usepackage{amsmath}

\title{MATH 1190 \& 1210 Meeting for Week 1}
\author{Josh}
\date{9 January 2026}

\begin{document}

\frame{\titlepage}

\frame{{Welcome!}

MATH 1210 has traditionally been a first course in calculus for students in business, social science, or some computational sciences. 

A few semesters ago, we began redesigning the course to better meet the needs of our partner departments. 

Some big content differences you may or may not know about:

\begin{itemize}
\item excluded content
    \begin{itemize}
    \item computing limits via algebraic manipulations
    \item problems involving the limit definition of the derivative
    \item implicit differentiation or related rates
    \item curve sketching
    \end{itemize}
\item included content
    \begin{itemize}
    \item introduction to summations
    \item introduction to probability, discrete random variables, continuous random variables, expectation, variance
    \end{itemize}
\end{itemize}
}

\frame{{Welcome!}



This semester, we will be fully using the growth-based assessment structure with some tweaks in the grading scale after some statistical analysis of previous semester.

\

}

\frame{{Welcome!}

\small

Assessment 1 targets
\begin{itemize}
\item new: F1, L1, L2,  D1
\end{itemize}
Assessment 2 targets 
\begin{itemize}
\item new: F2, D2, D3, A1
\item reassess: F1, L1, L2, D1
\end{itemize}
Assessment 3 targets 
\begin{itemize}
\item new: A2, A3, I1, I2
\item reassess: F2, D2, D3, A1
\end{itemize}
Assessment 4 targets 
\begin{itemize}
\item new: I3, I4, F3, P1, P2
\item reassess: A2,A3,I1, I2
\end{itemize}
Final Assessment targets
\begin{itemize}
\item FE targets: D1,D2,D3,I1,I3,I4,F3,L1,P1,P2
F3, L3, D1, D2, A2, A3, I1, I4, P1, P2
\item reassess: all others
\end{itemize}

}

\frame{{Welcome!}

Some changes have been made based on student and instructor feedback:
\begin{itemize}
\item We will combine Expectation, Variance, Continuous and Discrete Probability across two lessons
\item We will continue with the 5th scoring category called "Exceeds Expectations" has been added above "Proficiency"
\item Some scores of "Almost" will count toward students' grades
\item Students are now able to miss one core target and still be in the B range and two core targets and still be in the C range.
\item ``PAW" averages are capped at 100\% (in both growth-based and standard)
\end{itemize}

}

\frame{{Welcome!}

\vfill

{\bf That every instructor can communicate to all students how the grading scheme works is critical.}

\vfill

Please read the syllabus and the related pages on Canvas.

\vfill


\vfill

We will discuss how to score assessments in a growth-based course throughout the semester. What I want to leave you with regarding this is the following:

\

Proficient does not mean perfect!

\

Students and some instructors seemed to have this viewpoint in previous semesters and it's antithetical to the goals of growth-based.

}
\frame{{Welcome!}

Some additional notes about our meetings this semester:
\begin{itemize}
\item I have been working with Jessica Weaver-Kenney and Lori Hostetler from Learning Design & Technology to create a simplified course page on canvas that should help with student navigation and understanding.
\item In addition to our usual weekly meetings logistics, we will add some discussions on pedagogy and current research in the field
\item To this end, I will kindly ask each instructor to develop one active learning activity to be implemented in their classroom and then speak about how it went the following week.
\item Due to positive student feedback, we will continue with the application days and take a stab at revamping the gradient descent application to the one variable case

\end{itemize}

}

\frame{{Welcome!}
\LARGE Let's all have a great semester! Thanks in advance for all your hard work.}

\end{document}